Kurs:Algebraische Kurven (Osnabrück 2012)/Arbeitsblatt 30/latex
Kurs:Algebraische Kurven (Osnabrück 2012)/Arbeitsblatt 30/latex


\setcounter{section}{30}


  \zwischenueberschrift{Aufwärmaufgaben}


\inputaufgabe

{}

{


Zeige durch ein Beispiel, dass Lemma 30.2
nicht ohne die Voraussetzung gilt, dass der Körper
\definitionsverweis {algebraisch abgeschlossen}{}{} ist.


}

{} {}


\inputaufgabe

{}

{


Zeige, dass in den in
Beispiel 30.6
berechneten Schnittpunkten

\mathl{\neq (0,0)}{} der beiden Kurven ein
\definitionsverweis {transversaler Schnitt}{}{} vorliegt.


}

{} {}


\inputaufgabegibtloesung

{}

{


Es sei


\mavergleichskette

{\vergleichskette

{ K
}

{ = }{ {\mathbb C}
}

{  }{
}

{    }{
}

{   }{
}

}
{}{}{.}
Bestimme für die beiden affinen Kurven


\mathdisp {V \left(  Y-X^3  \right)  \text{ und } V \left(  Y^2-X^3  \right)} {  }


ihre Schnittpunkte zusammen mit den
\definitionsverweis {Schnittmultiplizitäten}{}{.}
Betrachte auch Schnittpunkte im ${\mathbb P}^{2}_{ {\mathbb C}}$ und bestätige
den Satz von Bézout
in diesem Beispiel.


}

{} {}


\inputaufgabegibtloesung

{}

{


Es sei


\mavergleichskette

{\vergleichskette

{ K
}

{ = }{ {\mathbb C}
}

{  }{
}

{    }{
}

{   }{
}

}
{}{}{.}
Wir betrachten die beiden ebenen algebraischen Kurven


\mathdisp {C=V \left(  X-Y^2  \right)  \text{ und } D=V \left(  Y^2-X^5  \right)} { . }


Bestimme die Schnittpunkte der beiden Kurven in der affinen Ebene und bestimme jeweils die Schnittmultiplizität. Bestimme auch die unendlich fernen Punkte der beiden Kurven
\zusatzklammer {also die zusätzlichen Punkte auf den projektiven Abschlüssen $\bar{C}$ und $\bar{D}$} {} {} und überprüfe damit die Schnittpunkte im Unendlichen. Bestätige abschließend, dass der Satz von Bézout in diesem Beispiel erfüllt ist.


}

{} {}


  \zwischenueberschrift{Aufgaben zum Abgeben}


\inputaufgabe

{6}

{


Es sei


\mavergleichskette

{\vergleichskette

{ C
}

{ \subseteq }{ {\mathbb P}^{2}_{K}
}

{  }{
}

{    }{
}

{   }{
}

}
{}{}{}
eine glatte Quadrik
\zusatzklammer {also eine Kurve vom Grad zwei} {} {}
über einem
\definitionsverweis {algebraisch abgeschlossenen Körper}{}{.}
Zeige, dass es eine Isomorphie der Kurve mit der projektiven Geraden ${\mathbb P}^{1}_{K}$ gibt.


}

{} {}


\inputaufgabe

{5}

{


Es sei $K$ ein
\definitionsverweis {algebraisch abgeschlossener Körper}{}{}
und


\mavergleichskette

{\vergleichskette

{C
}

{ \subset }{ {\mathbb P}^{2}_{K}
}

{  }{
}

{    }{
}

{   }{
}

}
{}{}{}
eine
\definitionsverweis {glatte Kurve}{}{}
vom Grad


\mavergleichskette

{\vergleichskette

{d
}

{ \geq }{ 2
}

{  }{
}

{    }{
}

{   }{
}

}
{}{}{.}
Zeige, dass es einen
\definitionsverweis {Morphismus}{}{}
\maabb {} { C } { {\mathbb P}^{1}_{K}
} {}
derart gibt, dass jede Faser aus maximal

\mathl{d-1}{} Punkten besteht.


}

{} {}


\inputaufgabe

{5}

{


Es sei


\mavergleichskette

{\vergleichskette

{ C
}

{ = }{ V_+ \left(  X^3+Y^3+Z^3  \right)
}

{ \subseteq }{ {\mathbb P}^{2}_{K}
}

{    }{
}

{   }{
}

}
{}{}{}
die \stichwort {Fermat-Kubik} {} über einem
\definitionsverweis {algebraisch abgeschlossenen Körper}{}{}
der
\definitionsverweis {Charakteristik}{}{}
$\neq 3$. Beschreibe explizit einen Morphismus
\maabb {} { C } { {\mathbb P}^{1}_{K}
} {,}
bei dem über jedem Punkt maximal zwei Punkte liegen.


}

{} {}


\inputaufgabe

{4}

{


Es sei


\mavergleichskette

{\vergleichskette

{ C
}

{ \subseteq }{ {\mathbb P}^{2}_{ {\mathbb C}}
}

{  }{
}

{    }{
}

{   }{
}

}
{}{}{}
der komplex-projektive Abschluss des Einheitskreises. Bestimme eine explizite bijektive Parametrisierung
\maabb {} { {\mathbb P}^{1}_{ {\mathbb C}} } { C
} {.}


}

{} {}


\inputaufgabe

{4}

{


Bestätige
den Satz von Bézout
für die beiden
\definitionsverweis {ebenen projektiven Kurven}{}{}


\mavergleichskette

{\vergleichskette

{ C
}

{ = }{ V_+ \left(  ZY^2-X^3  \right)
}

{  }{
}

{    }{
}

{   }{
}

}
{}{}{}
und


\mavergleichskette

{\vergleichskette

{ D
}

{ = }{ V_+ \left(  X^2+(Y-Z)^2-Z^2  \right)
}

{  }{
}

{    }{
}

{   }{
}

}
{}{}{.}


}

{Skizziere die Situation.} {}


\inputaufgabe

{4}

{


Bestätige
den Satz von Bézout
für die beiden
\definitionsverweis {ebenen projektiven Kurven}{}{}


\mavergleichskette

{\vergleichskette

{C
}

{ = }{ V_+ \left(  ZY-X^2  \right)
}

{  }{
}

{    }{
}

{   }{
}

}
{}{}{}
und


\mavergleichskette

{\vergleichskette

{D
}

{ = }{ V_+\left(  X^2+(Y-Z)^2-Z^2  \right)
}

{  }{
}

{    }{
}

{   }{
}

}
{}{}{.}


}

{Skizziere die Situation.} {}


\inputaufgabe

{4}

{


Bestätige
den Satz von Bézout
für die beiden monomialen Kurven, die affin durch


\mavergleichskette

{\vergleichskette

{ C
}

{ = }{ V(X^2-Y^3)
}

{  }{
}

{    }{
}

{   }{
}

}
{}{}{}
und


\mavergleichskette

{\vergleichskette

{ D
}

{ = }{ V(X^5-Y^4)
}

{  }{
}

{    }{
}

{   }{
}

}
{}{}{}
gegeben sind.


}

{} {}


\inputaufgabe

{5}

{


Es sei $R$ ein
\definitionsverweis {kommutativer Ring}{}{}
und seien $M, N$ $R$-Moduln. Ist
\maabbdisp {f} { M } { N
} {}
ein
$R$-\definitionsverweis {Modulhomomorphismus}{}{,}
so ist
\maabbeledisp {f^\ast} { \operatorname{Hom} \, (N,R) } { \operatorname{Hom} \, (M,R)
} {\varphi } {\varphi \circ f
} {,}
auch ein $R$-Modulhomomorphismus.


Es sei nun $0 \to M \to N \to P \to 0$ eine
\definitionsverweis {kurze exakte Sequenz}{}{}
von $R$-Moduln.


Zeige, dass dann die induzierte Sequenz


\mathdisp {0 \longrightarrow \operatorname{Hom} \, (P,R) \longrightarrow \operatorname{Hom} \, (N,R) \longrightarrow \operatorname{Hom} \, (M,R)} {  }


exakt ist. Man gebe auch ein Beispiel mit $R = \Z$, das zeigt, dass der letzte Pfeil im Allgemeinen nicht surjektiv ist.


}

{} {}


<< | Kurs:Algebraische Kurven (Osnabrück 2012) | >>


PDF-Version dieses Arbeitsblattes


 Zur Vorlesung
(PDF)
